% !TEX root = ../monograph.tex
\chapter{Assembly and the scope of the construction}
\label{ch19:assembly}

The preceding chapters distinguish finite algebraic identities, analytic
estimates for the lattice transformation, and consequences for the
continuum theory. This chapter states how they fit together. The construction and finite-order operator products use the estimates
proved in the preceding chapters. The sharper trace asymptotic is
separated below because it requires the additional current estimates
of Chapter~\ref{ch17:trace}.

\section{One family and the order of its limits}
\label{ch19:family}

Choose the compact matrix group $G$ with simple nonabelian real Lie
algebra, its trace inner product as in Section~\ref{ch01:gauge}, and a
positive soft-block penalty $b$. Fix a sufficiently large coarse inverse
coupling $\bar\beta$. The common schedule of
Chapter~\ref{ch10:trajectory} selects all block lengths, derivative
orders, support weights, coupling thresholds and admissible physical
periods $R_p$. At cutoff $N$ and period index $p$, shooting from the
Wilson measure gives its bare coupling $\beta_{N,p,N}$ with prescribed
coarse endpoint $\bar\beta$.

The same finite histories must define the coefficient normalizations
of every local insertion. One cannot separately choose convenient
histories for convergence, for free separating families and for the trace law.
The microscopic sections are fixed before $N,p$, and their comparison
maps act simultaneously on all labels and spectator correlation pairings. Any additional
finite stage requirements, such as the equation calibration floor,
belong to the selection of this one schedule.

For each fixed physical volume the first limit is the full limit
$N\to\infty$. Only afterwards is a subsequence of physical volumes
extracted. Given any increasing sequence of period indices, a single
refinement supplies all arities, degree bounds, label basis elements and
derivatives of the separated kernels. Auxiliary coefficient limit points
may be selected on ordered tails; they do not change this full inner
limit of the correlations. Scaling comparison sequences and the choice
of a phase are later operations on the same continuum family.

\section{Assembly of the construction}
\label{ch19:theorem-section}

The hypotheses used along the construction have now been supplied for
one common family. Chapters~\ref{ch09:renormalization} and
\ref{ch10:trajectory} give the actual block estimates and fixed-input
convergence. Chapter~\ref{ch11:insertions} supplies the marked comparisons,
including restarted histories and arbitrary fixed numbers of insertions.
Chapter~\ref{ch13:continuum} constructs the bounded approximations and
local product estimates. The free comparison and annular estimates in
Chapter~\ref{ch14:ward}, and the repeated estimate for a separating family in
Chapter~\ref{ch15:phases}, use these same objects. Thus they are not
additional assumptions in the following assembly statement.

\begin{theorem}[Assembly of the common Wilson family]
\label{ch19:conditional}
The chosen Wilson family has the following
properties. For every prescribed increasing volume sequence there is
one subsequence along which the cutoff limits converge to a hierarchy $S$ of smooth separated Euclidean
correlations and compatible finite-dimensional field spaces $F_d$.
Their dimensions and orthogonal representations agree with those of
the invariant free Maxwell comparison spaces. The filtration $(F_d)$ contains the unit field. Its inclusion maps are
injective, and its four coordinate derivatives commute.

There is a compact metrizable phase space $\Xi$, a full-support
probability measure $\nu$, and phase hierarchies $S^\xi$ with
$S=\int S^\xi\,d\nu$. All phases have the same field-null spaces,
Euclidean covariance and reflection positivity. Each phase reconstructs
a local relativistic quantum field theory with tempered Wightman
correlations, a faithful common field filtration, a positive-cone
translation spectrum and a unique invariant vacuum. Its weak field
commutant is scalar. These statements are uniform over the phases in
the finite-list senses specified in the preceding chapters.

The common real stress tensor is symmetric and conserved, with its
normalization fixed by local translation and rotation Ward identities.
Its relativistic continuation gives the local charge law
\eqref{ch16:charge-law}. The fixed OPE prescription of
Section~\ref{ch18:prescription} gives the phase-uniform finite-order
coefficient and remainder bounds \eqref{ch18:ope-bound}. Independently,
the fixed subtraction prescription
of Section~\ref{ch18:extension} extends all separated correlations to
finite-order distributions across collision diagonals, compatibly
with the field inclusions, translations, rotations and permutations.
\end{theorem}
\begin{proof}
Chapter~\ref{ch12:ultraviolet} takes the full cutoff
limit at each fixed period. The diagonal subsequence in the volume index works for every fixed
number of insertions and every scaling dimension. The two rank bounds
in Chapter~\ref{ch13:continuum} identify the dimensions of the resulting
field spaces, and its derivative maps give their differential calculus. That chapter also proves the repeated-word bound,
so bounded slab insertions in Chapter~\ref{ch14:ward} do not depend on
later reconstruction or operator products.

Coordinate reflection positivity, the free comparison and the annular
estimates give the conserved tensor, its local charges, Euclidean
rotations and Ward normalization in Chapter~\ref{ch14:ward}.
Chapter~\ref{ch15:phases} then compresses whole words to the mixed
ground space. The real approximation and moments for separating families prove
symmetries and the equality of field spaces in every character, beyond
what abstract commutative algebra alone gives.

The slab bounds supply joint boundary distributions in
Chapter~\ref{ch16:reconstruction}. Covariance gives the spectral cone,
analytic continuation gives locality, and word reconstruction gives
the relativistic fields. The unitary equivalence of the reconstructed and reflection-positive Hilbert spaces and the
phase-uniform low-energy estimate remove exceptional vacua. The same
chapter proves weak irreducibility and continues the stress charge law.

The finite OPE and contact extension follow from their separate arguments in
Chapter~\ref{ch18:products}. Each step preserves the originally chosen
family and the already selected correlation limits.
\end{proof}

\begin{corollary}[The trace law]
\label{ch19:conditional-trace}
In the phase $\xi_*$ of Section~\ref{ch15:selection}, let
$\Theta=\sum_aT_{aa}-S^{\xi_*}(\sum_aT_{aa})1$ be the centered trace and
$G_\Theta(r)=S_2^{\xi_*}(\Theta,\Theta;re_4,0)$ its two-point function.
Then
\[
 \lim_{r\downarrow0}r^8\log^2(1/r)G_\Theta(r)
                    =\frac{3\dim_\RR\gLie}{\pi^4}.
\]
Consequently the centered trace cannot agree near zero with any finite linear combination of products of the
mass-shell Laplace kernels specified in Section~\ref{ch17:exclusion}.
\end{corollary}
\begin{proof}
Apply the rate-sensitive tensor cancellation and prescribed-radius
squeeze of Chapter~\ref{ch17:trace}. The positive spectral measure
required there already follows from reconstruction.
\end{proof}

\section{Conclusions that are not supplied}
\label{ch19:scope}

The logarithmic ultraviolet law and the low-energy estimate describe
different regimes. The latter bounds spectral mass near zero by a
power; it does not show that the spectrum avoids an interval above the
vacuum. A positive mass gap is not proved.

Vacuum uniqueness within each phase does not imply uniqueness of the
phase, a point-mass phase measure, or uniqueness of the infinite-volume
limit. The construction uses a refinement of a prescribed volume
sequence. It does not prove a simultaneous continuum and infinite-volume
limit, independence of the block penalty or all normalization choices,
or a separate lattice realization of each character phase.

The trace comparison excludes the finite class explicitly defined in
Chapter~\ref{ch17:trace}. It does not rule out every generalized free
field or every infinite Wick expansion on the basis of one two-point
function. The OPE is an expansion to each fixed finite degree; no
infinite series convergence has been asserted. The contact extension
has a specified subtraction convention and may have larger local order
than an optimized prescription. Positivity at contacts, contact Ward
identities, the fact that inserting the identity leaves correlations unchanged, and uniqueness of contact
terms require additional normalization conditions.

The trace argument uses more than convergence to the free tensor. Its
low-moment estimate distinguishes field degree from scaling dimension:
a quadratic expression with extra spatial differences is still quadratic
in the field. Exact polynomial preservation by the linear response,
together with the normalized Gaussian fluctuation estimate, supplies the
required decay in Proposition~\ref{ch17:spatial-moments}. This is why
that additional calculation is needed after reconstruction.
